\section{A matrix-cone homomorphism with a full corner}
\label{sec:tensor}

We first construct a full square-zero element in a finite tensor power of
a unital \(C^*\)-algebra without characters. Its polar decomposition will
then give the matrix-cone homomorphism needed in the next section.
The absence of characters is detected by finitely many commutators.
The main issue is to combine those finitely many tests into one element
that survives every irreducible representation of a maximal tensor
product, including representations that are not spatial tensor products.

\subsection{Algebraic independence of commuting factors}

\begin{lemma}
\label{lem:tensor-independence}
Let \(M\subseteq B(H)\) be a von Neumann factor with \(1_M=1_H\), and let
\(B\subseteq M'\) be a unital algebra. The multiplication map
\[
 M\odot B\longrightarrow B(H),\qquad a\otimes b\longmapsto ab,
\]
is injective. Consequently, multiplication is injective on the complex
algebraic tensor product of any finite family of commuting von Neumann
factors on the same Hilbert space, all with unit \(1_H\).
\end{lemma}

\begin{proof}
Fix complex linearly independent elements \(b_1,\ldots,b_l\in B\).
It suffices to show that
\[
 S=\left\{(a_1,\ldots,a_l)\in M^{1\times l}:
                  \sum_{i=1}^l a_i b_i=0\right\}
\]
is zero. This row space is weak-operator closed and invariant under
diagonal multiplication by \(M\) on either side. The right invariance
uses \(B\subseteq M'\). Let \(J\subseteq M_l(M)\) consist of the matrices
whose rows belong to \(S\). It is a weak-operator closed left ideal:
each row of a left matrix product is an \(M\)-linear combination of the
original rows.

We give the support-projection argument for the form of this ideal.
For \(x\in J\), we have \(x^*x\in J\). Functional calculus with functions
vanishing at zero, followed by weak-operator closure, puts the support
projection of \(x^*x\) in \(J\): one may first take
\(x^*x(x^*x+\varepsilon1)^{-1}\) and then let
\(\varepsilon\downarrow0\). Finite joins of these projections also
belong to \(J\), since the join of finitely many projections is the
support of their sum. Their increasing join \(p\) therefore belongs to
\(J\); this is a net of finite joins, so no separability assumption is
involved. Every \(x\in J\) satisfies \(x=xp\), and \(p\in J\) gives
\[
                         J=M_l(M)p.
\]

For any unitary \(u\in M\), put \(U=\operatorname{diag}(u,\ldots,u)\).
Right invariance and its counterpart for \(u^*\) give \(JU=J\).
The right support projection of \(xU\) is the conjugate by \(U\) of the
right support projection of \(x\). Taking their joins gives
\[
                             U^*pU=p.
\]
Thus every entry of \(p\) belongs to the center of \(M\), and is scalar.
Each row of \(p\) is consequently a scalar relation among the independent
\(b_i\), so \(p=0\). Hence \(J=0\) and \(S=0\), proving the first assertion.

For finitely many commuting factors, apply the first assertion to one
factor and the algebra generated by the remaining factors, and induct
on their number.
\end{proof}

\begin{lemma}
\label{lem:tensor-grouped-factors}
Let \(\pi:B_1\otmax\cdots\otmax B_q\to B(H)\) be an irreducible
representation, where the \(B_j\) are nonzero unital
\(C^*\)-algebras. The weak-operator closures of the represented tensor
legs are commuting von Neumann factors. The same is true after grouping
the legs according to any partition into nonempty blocks.
\end{lemma}

\begin{proof}
Every represented block contains \(\pi(1)=1_H\). The represented blocks
commute, and this commutation persists under
weak-operator closure. A central element of the closure of one block
commutes both with that block and with every complementary block.
It therefore belongs to the commutant of the full irreducible image,
which consists of the scalars. This proves factorhood for each block.
\end{proof}

These lemmas concern algebraic tensor independence. They require neither
faithfulness of the restrictions to the tensor legs nor a spatial
decomposition of the represented factors.

\subsection{Labels that remain independent in representations}

\begin{lemma}
\label{lem:tensor-detecting-space}
Let \(D\) be a nonzero unital \(C^*\)-algebra without characters.
There is a finite-dimensional real subspace \(V\subseteq D_{\sa}\)
such that the image of \(V\) contains a noncommuting pair in every
nonzero unital representation of \(D\).
\end{lemma}

\begin{proof}
The closed ideal generated by the commutators of selfadjoint elements
is all of \(D\). Otherwise the quotient by this ideal would be a nonzero
unital commutative \(C^*\)-algebra, and a character of the quotient would
give a character of \(D\).

Choose a finite algebraic expression in these commutators within distance
less than one of \(1_D\). This expression is invertible, so the finitely
many commutators it uses already generate the unit ideal. Let \(V\) be
the real span of their selfadjoint entries. If a unital representation
sent \(V\) to a commuting space, it would annihilate all these
commutators, and hence their unit ideal. This is impossible.
\end{proof}

\begin{lemma}
\label{lem:tensor-real-labels}
Let \(V\) be a finite-dimensional real vector space with
\(m=\dim_{\R}V\ge2\). For every integer \(r\ge1\), there are \(n\ge1\)
and tensors \(x_1,\ldots,x_r\in V^{\otimes_{\R}n}\) with the following
property: whenever \(P_j:V\to E_j\) are real linear maps of rank at least
two, the tensors
\[
               (P_1\otimes\cdots\otimes P_n)(x_i),
               \qquad 1\le i\le r,
\]
are linearly independent over \(\R\).
\end{lemma}

\begin{proof}
It suffices to consider surjections \(P_j:V\to\R^2\).
Indeed, a map of rank at least two can be postcomposed, on its
finite-dimensional image, with a surjection onto \(\R^2\); a dependence
before this postcomposition would remain a dependence afterward.

Choose \(n\) such that
\[
                         2^n>2mn+r-1,
\]
and put \(W=V^{\otimes_{\R}n}\). Let
\(\mathcal P=\operatorname{Surj}(V,\R^2)^n\), an open subset of a real
vector space of dimension \(2mn\). Consider the smooth map
\[
 \begin{split}
 \Phi:\mathcal P\times S^{r-1}\times W^r
     &\longrightarrow (\R^2)^{\otimes_{\R}n},\\
 (P,u,x)&\longmapsto
 (P_1\otimes\cdots\otimes P_n)\left(\sum_{i=1}^r u_i x_i\right).
 \end{split}
\]
Its differential in the \(x\) variables is surjective everywhere:
some \(u_i\ne0\), and the tensor product of the \(P_j\) is surjective.
Thus \(\Phi^{-1}(0)\) is a smooth manifold of dimension
\[
                 2mn+r-1+rm^n-2^n<rm^n=\dim_{\R}W^r.
\]
By Sard's theorem, its projection to \(W^r\) has Lebesgue measure zero.
The manifolds here are second countable, so this assertion can be
applied on a countable atlas; no compactness of \(\mathcal P\) is needed.
Choose \(x\in W^r\) outside that projection. For this one tuple, no
nonzero real dependence occurs for any tuple of surjections: any
nonzero coefficient vector can be normalized to lie in \(S^{r-1}\).
\end{proof}

\begin{corollary}
\label{cor:tensor-represented-labels}
Let \(D\) and \(V\) be as in Lemma~\ref{lem:tensor-detecting-space}.
For each \(r\ge1\), there are \(n\ge1\) and selfadjoint elements
\(x_1,\ldots,x_r\in D^{\otimes n}\) whose images are complex linearly
independent under every irreducible representation of
\(D^{\otimes n}\otmax B\), for every nonzero unital \(C^*\)-algebra \(B\).
Here the elements \(x_i\) occupy the first \(n\) tensor legs.
\end{corollary}

\begin{proof}
Take the real tensors supplied by Lemma~\ref{lem:tensor-real-labels}
and map them naturally into \(D^{\otimes n}\). They are selfadjoint,
since they are real linear combinations of elementary tensors with
selfadjoint entries.

Fix such an irreducible representation and write \(\pi_j\) for its
restriction to the \(j\)-th copy of \(D\). Each \(\pi_j\) is unital, so
Lemma~\ref{lem:tensor-detecting-space} implies
\(\dim_{\R}\pi_j(V)\ge2\). Apply
Lemma~\ref{lem:tensor-real-labels} to the real maps
\(V\to\pi_j(V)\).

The resulting independence also holds after complexification.
Indeed, a real basis of a space of selfadjoint operators is complex
linearly independent: taking adjoints of a complex relation separates
its real and imaginary parts. Hence the complexification of each
\(\pi_j(V)\) embeds in the represented algebra. Tensoring these
injections shows that the real independent tensors remain complex
independent in the complex algebraic tensor product of the represented
legs: the latter tensor product contains the complexification of the
real tensor product injectively. By
Lemma~\ref{lem:tensor-grouped-factors} their weak closures are
commuting factors, and Lemma~\ref{lem:tensor-independence} preserves this
independence under multiplication.
\end{proof}

\subsection{A full square-zero element}

The existence of two full orthogonal positive elements in a finite
maximal tensor power is already established in
\cite[Theorem~1.4]{Fixed231}. The bridge construction below is related to
Kirchberg--R\o rdam's proof of \cite[Proposition~6.1, (i)$\Rightarrow$(v)]{KR},
which combines finitely many bridges with orthogonal full tensor labels.
Here we use representation-stable linear independence and give the full
construction of the square-zero element needed for the cone map.

The labels allow linear independence to replace orthogonality when
combining finitely many algebraic witnesses. We first construct two full
orthogonal positive elements, and then use a second label family to
retain all bridges between them in a single square-zero element.

\begin{lemma}
\label{lem:tensor-full-square-zero}
If \(D\) is a nonzero unital \(C^*\)-algebra without characters, then
some finite maximal tensor power \(P\) of \(D\) contains a full element
\(w\) satisfying \(w^2=0\) and \(\|w\|=1\).
\end{lemma}

\begin{proof}
Choose \(V\) as in Lemma~\ref{lem:tensor-detecting-space}, and let
\(v_1,\ldots,v_m\) be a real basis of \(V\). Take the labels
\(x_1,\ldots,x_m\in D^{\otimes n}\) from
Corollary~\ref{cor:tensor-represented-labels}. In
\[
                   P_0=D\otmax D^{\otimes n}\otmax D^{\otimes n}
\]
define selfadjoint elements
\[
 s=\sum_{i=1}^m v_i\otimes x_i\otimes1,
 \qquad
 t=\sum_{i=1}^m v_i\otimes1\otimes x_i.
\]
We first show that their images do not commute in any irreducible
representation \(\pi\) of \(P_0\). Denote the restrictions to the three
displayed blocks by \(\pi_0,\pi_L,\pi_R\). Their weak closures are
commuting factors by Lemma~\ref{lem:tensor-grouped-factors}, and each
of the two label families has independent images by
Corollary~\ref{cor:tensor-represented-labels}. The expansion
\[
 \pi([s,t])=
 \sum_{i,j=1}^m
       \pi_0([v_i,v_j])\,\pi_L(x_i)\,\pi_R(x_j)
\]
and Lemma~\ref{lem:tensor-independence} show that
\(\pi([s,t])=0\) would imply
\(\pi_0([v_i,v_j])=0\) for all \(i,j\).
This contradicts the defining property of \(V\) for the unital
representation \(\pi_0\). Independence of the images of the \(v_i\)
themselves is not required.

Set \(z=[s,[s,t]]\), which is selfadjoint. We claim that both \(z_+\)
and \(z_-\) are full in \(P_0\). Fix an irreducible representation
\(\pi\), and suppose first that \(\pi(z)\ge0\). For every vector
\(\xi\), the real function
\[
 f_\xi(\lambda)=
 \left\langle e^{i\lambda\pi(s)}\pi(t)e^{-i\lambda\pi(s)}
                       \xi,\xi\right\rangle,\qquad \lambda\in\R,
\]
is bounded and twice continuously differentiable, with
\[
 f_\xi''(\lambda)=
 -\left\langle e^{i\lambda\pi(s)}\pi(z)e^{-i\lambda\pi(s)}
                       \xi,\xi\right\rangle\le0.
\]
A bounded concave function on the whole real line is constant.
Consequently
\[
                f_\xi'(0)=\langle i\pi([s,t])\xi,\xi\rangle=0
                \qquad\text{for every }\xi.
\]
The operator \(i\pi([s,t])\) is selfadjoint, so it must be zero, a
contradiction. The assumption \(\pi(z)\le0\) leads to the same
contradiction using bounded convex functions. Thus both the positive
and negative parts of \(\pi(z)\) are nonzero. Functional calculus commutes with \(\pi\),
so \(z_+\) and \(z_-\) survive every irreducible representation and
are full.

Let \(J\) be the closed ideal generated by \(z_-P_0z_+\).
The algebraic two-sided ideals
\[
 I_\pm=\operatorname{span}\{a z_\pm b:a,b\in P_0\}
\]
are both norm dense in \(P_0\). Their products lie in \(J\), since
\[
             (a z_- b)(c z_+ d)=a(z_-bc z_+)d.
\]
Approximating \(1_{P_0}\) from both \(I_-\) and \(I_+\) shows that
\(1_{P_0}\in J\), and hence \(J=P_0\). An algebraic expression
within distance less than one of \(1_{P_0}\) uses only finitely many
elements of \(z_-P_0z_+\) and their adjoints. We can therefore choose
\[
                   c_i=z_-d_i z_+,\qquad 1\le i\le r,
\]
whose closed ideal is \(P_0\). Notice that \(c_i c_j=0\) for all
\(i,j\), because \(z_+z_-=0\).

Choose new labels \(x_1',\ldots,x_r'\in D^{\otimes n'}\) from
Corollary~\ref{cor:tensor-represented-labels}, and put
\[
 P=P_0\otmax D^{\otimes n'},\qquad
                       w_0=\sum_{i=1}^r c_i\otimes x_i'.
\]
The identities \(c_i c_j=0\) give \(w_0^2=0\).
To prove fullness, take any irreducible representation of \(P\),
with restrictions \(\rho_0\) to \(P_0\) and \(\rho_1\) to the new
label block. The images \(\rho_1(x_i')\) are independent by
Corollary~\ref{cor:tensor-represented-labels}, allowing \(P_0\) as
the additional algebra. The two grouped weak closures are factors
by Lemma~\ref{lem:tensor-grouped-factors}. Thus
Lemma~\ref{lem:tensor-independence} implies
\[
       \sum_i\rho_0(c_i)\rho_1(x_i')=0
       \quad\Longrightarrow\quad
       \rho_0(c_i)=0\quad\text{for all }i.
\]
The latter conclusion is impossible, since the \(c_i\) generate the
unit ideal and \(\rho_0\) is unital. Hence \(w_0\) is full.
In particular \(w_0\ne0\), and \(w=w_0/\|w_0\|\) has the required
properties. The exponents \(n,n'\), and hence the exponent of \(P\),
depend on the finite detecting space and the finite bridge family;
no bound uniform over \(D\) is asserted.
\end{proof}

\subsection{The cone homomorphism}

Let
\[
 C=\{f\in C([0,1],M_2):f(0)\in\C1_{M_2}\},
 \qquad \iota(t)=t.
\]
Thus \(C\) is the unitization of \(C_0((0,1],M_2)\).

\begin{proposition}
\label{prop:full-cone}
For every nonzero unital \(C^*\)-algebra \(D\) without characters,
there are an integer \(m\ge1\), the maximal tensor power
\(P=D^{\otimes m}\), and a unital \(*\)-homomorphism
\(\eta:C\to P\) such that \(\eta(\iota e_{11})\) is full in \(P\).
\end{proposition}

\begin{proof}
Let \(P\) and \(w\) be supplied by
Lemma~\ref{lem:tensor-full-square-zero}. Take the polar decomposition
\(w=u|w|\) in \(P^{**}\), and put \(p=u^*u\), \(q=uu^*\).
Since \(w^2=0\), the closure of the range of \(w\) lies in its kernel;
hence \(q\le1-p\), since \(p\) is the projection onto the orthogonal
complement of that kernel in a faithful normal representation of \(P^{**}\).
The elements
\[
           E_{11}=p,\qquad E_{22}=q,\qquad
           E_{21}=u,\qquad E_{12}=u^*
\]
are therefore matrix units with unit \(p+q\).

The positive element \(h=|w|+|w^*|\) has norm one, support \(p+q\),
and commutes with these matrix units. Indeed, the two summands have
orthogonal supports, and the polar decomposition gives
\[
 hE_{11}=|w|,\qquad hE_{22}=|w^*|,\qquad
 hE_{21}=w,\qquad hE_{12}=w^*.
\]
Commuting functional calculus now defines a \(*\)-homomorphism
\[
 \begin{split}
 \varphi:C_0((0,1],M_2)&\longrightarrow P^{**},\\
 f\otimes e_{ij}&\longmapsto f(h)E_{ij},
       \qquad f\in C_0((0,1]).
 \end{split}
\]
It takes its values in \(P\). To see this directly, approximate \(f\)
uniformly by polynomials vanishing at zero. Each corresponding
term \(h^kE_{ij}\), \(k\ge1\), belongs to \(P\), because
\(hE_{ij}\in P\). Thus \(f(h)E_{ij}\in P\).
Unitizing \(\varphi\) gives a unital homomorphism \(\eta:C\to P\),
and
\[
                  \eta(\iota e_{21})=w,\qquad
                  \eta(\iota e_{11})=|w|.
\]
No assertion that \(u\) belongs to \(P\), or that \(p+q=1_P\), is
needed for this unitization.

Finally, \(w\) and \(|w|\) generate the same closed ideal.
The ideal of \(w\) contains \(w^*w\) and hence \(|w|\) by functional
calculus. Conversely,
\[
          w(|w|+\varepsilon1_P)^{-1}|w|\longrightarrow w
          \qquad(\varepsilon\downarrow0)
\]
in norm, so the ideal of \(|w|\) contains \(w\).
Fullness of \(w\) therefore implies fullness of
\(\eta(\iota e_{11})=|w|\).
\end{proof}
